\section{Обучение и оценка качества}

\subsection{Обучающая разметка}

Для каждой позиции обучающей выборки сначала выполняются все \(K\) внутренних шагов. Затем вычисляются потери \(L_{t,0},\ldots,L_{t,K}\) и оракульное время \(\tau_t^\star\) по формуле~\eqref{eq:oracle_stopping_time}. После этого обучается модель распределения \(\widehat p_{t,k}\).

Если используется точечная разметка, целевая функция имеет вид
\begin{equation}
  \mathcal L_{\mathrm{point}}
  =-\sum_t\sum_{k=0}^{K}\log \widehat p_{t,k}(\tau_t^\star).
  \label{eq:point_loss}
\end{equation}
Если используется сглаженная разметка, применяется функция~\eqref{eq:distribution_loss}. Оба варианта нужно сравнить, поскольку точечная разметка проще, а сглаженная лучше отражает близкие по риску шаги.

\subsection{Качество языкового прогноза}

Основная метрика качества прогноза --- средняя отрицательная логарифмическая потеря при фактическом времени остановки:
\begin{equation}
  \overline L=\frac{1}{N}\sum_{t=1}^{N}-\log Q_{t,\tau_t}(X_{t+1}).
  \label{eq:mean_nll}
\end{equation}
Дополнительно фиксируется доля совпадений верхнего варианта с правильной языковой единицей:
\begin{equation}
  \mathrm{Acc}=\frac{1}{N}\sum_{t=1}^{N}
  \ind\left\{\argmax_{a\in\A}Q_{t,\tau_t}(a)=X_{t+1}\right\}.
  \label{eq:accuracy}
\end{equation}
Эти метрики нужны вместе: логарифмическая потеря чувствительна к качеству распределения, а доля совпадений показывает поведение верхнего варианта.

\subsection{Качество предсказания времени}

Для самой модели времени используются три показателя. Первый --- отрицательная логарифмическая потеря по \(\tau^\star\):
\begin{equation}
  \mathrm{NLL}_{\tau}=-\frac{1}{N(K+1)}\sum_t\sum_{k=0}^{K}\log\widehat p_{t,k}(\tau_t^\star).
  \label{eq:tau_nll}
\end{equation}
Второй --- средняя абсолютная ошибка математического ожидания:
\begin{equation}
  \mathrm{MAE}_{\tau}=\frac{1}{N(K+1)}\sum_t\sum_{k=0}^{K}
  \left|\sum_{j=0}^{K}j\widehat p_{t,k}(j)-\tau_t^\star\right|.
  \label{eq:tau_mae}
\end{equation}
Третий --- ошибка калибровки для события \(\tau^\star\le k\). Она показывает, совпадает ли заявленная вероятность уже наступившего оптимума с фактической частотой.

\subsection{Вычислительная стоимость}

Стоимость измеряется средним числом внутренних шагов
\begin{equation}
  \overline\tau=\frac{1}{N}\sum_{t=1}^{N}\tau_t
  \label{eq:mean_tau}
\end{equation}
и долей экономии относительно полного запуска:
\begin{equation}
  \mathrm{Save}=1-\frac{\overline\tau}{K}.
  \label{eq:saving}
\end{equation}
Если учитывается начальный шаг \(0\), в знаменателе можно использовать \(K+1\); в экспериментах нужно зафиксировать один вариант и не менять его между таблицами.

